\documentclass{resume} % Use the custom resume.cls style

\usepackage[left=0.4 in,top=0.4in,right=0.4 in,bottom=0.4in]{geometry} % Document margins
\usepackage{changepage}
\usepackage{fontawesome}
\usepackage{ragged2e}
\usepackage[english,ngerman]{babel}
\newcommand{\tab}[1]{\hspace{.2667\textwidth}\rlap{#1}} 
\newcommand{\itab}[1]{\hspace{0em}\rlap{#1}}
\newcommand{\linkedin}{\faLinkedin\hspace{0.2em}}
\newcommand{\github}{\faGithub\hspace{0.2em}}
\name{Jens Ettl} % Your name
\address{+49 176 78867038 \\ Karlsruhe, Deutschland \\ \today} 
\address{\href{mailto:ettljens@gmail.com}{ettljens@gmail.com} \\ \linkedin \href{https://www.linkedin.com/in/jens-ettl-807578211/}{www.linkedin.com/in/jens-ettl} \\ \github \href{https://github.com/xy}{github.com/jensettl} \\ \href{www.jensettl.com}{www.jensettl.com}
} 


\begin{document}

% https://www.instagram.com/reel/C4yVB9Lo_HP/
%
% Update some Sections and add some new ones
%

% \begin{rSection}{Über mich}
%     % \raggedright{Als Student der Wirtschaftsinformatik mit über einjähriger Erfahrung als Product Owner bringe ich umfassende Kenntnisse in der Konzeption und Entwicklung von Apps und MVPs mit. Meine Fähigkeiten umfassen die Zusammenarbeit mit Kunden, die Erstellung präziser User Stories und Akzeptanzkriterien sowie die erfolgreiche Anwendung agiler Methoden wie SCRUM}
%     \raggedright{
%         Ich bin Jens Ettl, ein Student der Wirtschaftsinformatik mit über einjähriger Erfahrung in der Entwicklung 
%         und Konzeption von Produkten nach Scrum-Methoden in der Rolle des Product Owners. 
%         Meine Motivation liegt darin, durch meine Arbeit Mehrwerte zu schaffen und einen positiven Beitrag in der IT-Welt zu leisten.

%     }
    
% \end{rSection}

\begin{rSection}
    {Wichtigste Leistungen und Fähigkeiten}

    \begin{tabular}{ccc}
        \begin{minipage}{0.3\textwidth}
        \begin{itemize}
        \item \raggedright{Scrum-Expertise für agile Methoden}
        \item Kommunikationfähigkeit
        \end{itemize}
        \end{minipage} &
        \begin{minipage}{0.3\textwidth}
        \begin{itemize}
        \item \raggedright{Erfahrung mit ERP- und CRM-Systemen}
        \item Schnelle Auffassungsgabe
        \end{itemize}
        \end{minipage} &
        \begin{minipage}{0.3\textwidth}
        \begin{itemize}
        \item Stets auf dem neusten Stand der IT-Technologien
        \item Kreative Lösungsansätze
        \end{itemize}
        \end{minipage}
        \end{tabular}

\end{rSection}

%-----------------------------------------------------------------------------------------

\begin{rSection}{Berufserfahrung}

    \textbf{Productowner} \hfill Okt 2021 - März 2023\\
    CAS Software AG \hfill \textit{Karlsruhe, Deutschland}
    \begin{itemize}
        \itemsep -3pt {}
        \item \raggedright{Konzeption und Entwicklung von Apps und MVPs, sowie Umgang mit CRM- und ERP-Systemen.}
        \item \raggedright{Direktes Arbeiten mit Kundenanforderungen, inklusive Erstellung von User-Stories und Akzeptanzkriterien}
        \item \raggedright{Planung und Durchführung von Sprint-Reviews und Retrospektiven nach SCRUM.}
    \end{itemize}

    \textbf{Vorstand der Fachschaft Wirtschaftsinformatik} \hfill Okt 2020 - März 2022\\
    Technische Hochschule Karlsruhe \hfill \textit{Karlsruuhe, Deutschland}
    \begin{itemize}
        \itemsep -3pt {}
        \item \raggedright{Steigerung der Teamfähigkeit durch die Planung und Durchführung von Events und Workshops, sowie das Management von Konflikten zwischen Studierenden und Dozenten.} 
        \item \raggedright{Generierung neuer Firmenbeziehungen}
        \item \raggedright{Vertretung der Studierendeninteressen gegenüber der Hochschulleitung und Mitwirkung bei politischen Entscheidungen.}
        \item \raggedright{Kommunikation mit der Hochschulleitung, Erstellung von Präsentationen und Dokumentationen zur Unterstützung der Fachschaftsarbeit.}
    \end{itemize}
\end{rSection}

%----------------------------------------------------------------------------------------
% TECHINICAL STRENGTHS	
%----------------------------------------------------------------------------------------
\begin{rSection}{Fähigkeiten und Fertigkeiten}
    
    \begin{tabular}{ @{} >{\bfseries}l @{\hspace{6ex}} l }
        Technical Skills & Programmiersprachen (Python, JavaScript), UML, BPMN, Adobe XD,                            \\
        & SQL, Firebase, VisualStudioCode, Docker, Confluence, Jira, Microsoft Office               \\
        \\
        Soft Skills     & Teamwork, Kommunikation, Problemlösung, Strukturiertheit, Pünktlichkeit,                  \\
        & Kreativität, Schnelle Auffassungsgabe                                                     \\
        \\
        Sonstiges       & Englisch \textit{(Verhandlungssicher)}, Französisch \textit{(vertiefte Grundkenntnisse)}, \\
        & Führerschein A \& B
    \end{tabular}\\
\end{rSection}

\begin{rSection}{Bildungsweg}
    {\bf Hochschule Karlsruhe - Technik und Wirtschaft}, Bachelor Wirtschaftsinformatik  \hfill {Erwartet 2024}
    \begin{itemize}
        \item Tutor in den Modulen Datenbanken 1 und Statistik u. Business Intelligence.
        \item \raggedright{Wahlmodule: Erfolgreiche Durchführung von Data Science Projekten, Datenanalyse mit JMP, Machine Learning, User Centered Design}\\
    \end{itemize}

    {\bf Karlsruhe Institut für Technologie}, Bachelor Informatik  \hfill {\textit{Abgebrochen} 2016 - 2019}
\end{rSection}
%----------------------------------------------------------------------------------------
\begin{rSection}{Zusätzliche Aktivitäten}
    \begin{itemize}
        \item 	Teilnahme an Workshops zu SCRUM und Qualitätsmanagement.
        \item	Mentor an der Hochschule für Erstisemestler und Quereinsteiger.
        \item   Sprachaufenthalt in England 2014
    \end{itemize}


\end{rSection}
%----------------------------------------------------------------------------------------
%	EDUCATION SECTION

%----------------------------------------------------------------------------------------





\newpage


%----------------------------------------------------------------------------------------
%	WORK EXPERIENCE SECTION
%----------------------------------------------------------------------------------------
\begin{rSection}{Projekte}
    \vspace{0.5em}
    \justifying{In meinem Studium und meiner Freizeit  habe ich sowohl eigenständig als auch im Team eine Vielzahl von Projekten durchgeführt, 
    die ein breites Spektrum an Themen umfassen.
    Diese reichen von der Planung und Konzeption mithilfe von Office-Produkten wie Excel und Word 
    bis hin zur Anwendung von SCRUM-Methoden für die Entwicklung und Durchführung:}

    
    \begin{adjustwidth}{0.5cm}{}
        \vspace{-1.5em}
        \item \textbf{Portfolio Website} \textit{(React, Tailwind, Firebase Hosting)} \hfill {Januar 2024} \\
        \raggedright{Entwurf, Entwicklung und Verwaltung meiner persönlichen Online-Visitenkarte zur Präsentation meiner Projekte, Fähigkeiten und
        Erfahrungen. \href{https://www.jensettl.com//}{(Online verfügbar unter www.jensettl.com)}}

        \item \textbf{Identifizierung fehlerhafter Produkte} \textit{(Machine Learning, Pandas, Python)} \hfill {März 2023} \\
        \justifying{Aufbereitung und Analyse eines Datensatzes aus der Stahlproduktion mit Schwerpunkt auf Oberflächenunregel-mäßigkeiten.
        Implementierung datengesteuerter Lösungen zur frühzeitigen Erkennung fehlerhafter Stahlprodukte
        unter Anwendung von Algorithmen des maschinellen Lernens.}
    
        \item \textbf{Supply Chain Management Simulation} \textit{(Typescript, FastAPI, SCM, Excel)}  \hfill {Dezember 2022 - März 2023} \\
        \justifying{Durchführung eines Planspiels, um eine Lieferkette in der Fahrradindustrie zu simulieren. Hierfür wurde Microsoft Excel genutzt.
        Anschließende Entwicklung von Frontend und Backend mit Typescript und FastAPI um einen Planungszeitraum zu berechnen und zu simulieren. 
        Die Ergebnisse wurden in einer Abschlusspräsentation vorgestellt.}
    
        \item \textbf{Anwendungsprojekt SmartCampus} \textit{(Planung, Konzeption, Entwicklung)} \hfill {April 2022 - August 2022} \\
        \justifying{Evaluation und Implementierung sensorbasierter Mehrwertdienste auf dem Hochschulcampus. 
        Regelmäßiger Austausch mit dem Kunden um die Anforderungen zu verstehen und sicherzustellen, dass die Budgetplanung eingehalten wurde. 
        Durchlaufen der 4 Phasen des Lebenszyklus eines Projekts: Initiierung, Planung, Durchführung und Abschluss.}
        
        \item \textbf{Automatisierung eines Kundenauftragsprozesses} \textit{(BPMN, Java, Camunda)} \hfill {Februar 2022} \\
        \justifying{{Abbilden eines Kundenauftragsprozesses in BPMN und Untersuchung auf Automatisierungsmöglichkeiten.
        Anschließende Implementierung der Automatisierungsentwürfe als Gruppe und Simulation mit Camunda.}}
    
        \item \textbf{gRPC Client-Server Kommunikation} \textit{(JavaScript, gRPC)} \hfill {Oktober 2021} \\
        {Entwicklung einer effizienten Client-Server Kommunikation mit gRPC im Rahmen der Veranstaltung “Automatisierung von Geschäftsprozessen”.
        Dokumentation und Code wurde mit GitHub verwaltet. Aufgaben wurden mit Kanban in Trello organisiert.}
        
        \item \textbf{Git Workshop} \textit{(Git, Markdown)} \hfill {September 2021} \\
        {Entwicklung eines Git-Workshops für Studenten zum Erlernen der Grundlagen von Git und GitHub.
        Der Workshop fand an der Hochschule für Technik und Wirtschaft Karlsruhe statt.}
    
        \item \textbf{Untersuchung von Unternehmensdaten auf Compliance} \textit{(Datenanalyse, JMP)} \hfill {Mai 2021} \\
        {Untersuchung realer Unternehmensdaten mit Hilfe von statistischen Methoden, Entscheidungsbäumen und neuronalen Netzen,
        welche Arbeitsschritte in einem Produktionsprozess die Einhaltung der Vorschriften beeinflussen.}
    
        \item \textbf{Automatisierungen im Lebensalltag} \textit{(Python, Hosting)}  \hfill {Februar 2024} \\
        {Konzeption und Implementierung von Automatisierungen um meinen Lebensalltag zu erleichtern. Darunter ein automatisiertes Aufräumsystem und ein Reporting-Tool für meine Finanzen.} 
        
        \item \textbf{Modellierung eines IT-Modells} \textit{(Konzeption, Use-Cases, Persona)} \hfill {März 2020 - Juni 2020} \\
        {In der Veranstaltung "Modellierung von IT-Systemen" ging es um die Konzeption und Planung einer App mit Fokus auf den Kundenanforderungen, Use-Cases bis hin zur Abbildung der Prozesse mit BPMN und UML.}
    
    \end{adjustwidth}
    
\end{rSection}


%----------------------------------------------------------------------------------------
\begin{rSection}{Auszeichnungen und Zertifikate}
    \vspace{0.5em}

    \begin{itemize}
        \item \textbf{Arbeitszeugnis}, CAS Software AG
        \item \textbf{Modern JavaScript}, Udemy
        \item \textbf{Grundlagen des Onlinemarketings}, Google Zukunftswerkstatt
        \item \textbf{RPA Starter}, UI Path
    \end{itemize}


\end{rSection}

\begin{rSection}{Freizeit und Hobbies}
    \vspace{0.5em}
    \begin{itemize}
        \item \raggedright{\textbf{Sport und Aktivitäten:} Ich genieße Aktivitäten wie Wandern, Badminton, Klettern, Motorradfahren, Brettspiele, Schach, Skifahren und Fitness im Studio. Außerdem probiere ich gerne neue Sportarten aus.}
        \item \raggedright{\textbf{Networking und Events:} Ich nehme regelmäßig an Networking-Veranstaltungen im Raum Karlsruhe teil, darunter Start-Up Gründertreffen und verschiedene Meetups.}
        \item \raggedright{\textbf{Kulinarisches Experimentieren:} Ich liebe es, neue Rezepte auszuprobieren und mich kreativ auszuleben.} 
        \item \raggedright{\textbf{Fortbildung:} Mein Interesse an lebenslangem Lernen zeigt sich durch das Lesen von Fachliteratur und Romanen sowie das Anhören von Podcasts und TED-Talks.}
    \end{itemize}


\end{rSection}



\end{document}
